% !TEX root = ../main.tex
% Sources: Gaberdiel, Persson, Ronellenfitsch and Volpato, arXiv:1211.7074;
% Persson and Volpato, arXiv:1312.0622; Dijkgraaf, Pasquier and Roche 1990;
% Gaberdiel, Hohenegger and Volpato, arXiv:1106.4315; Gritsenko and Clery,
% arXiv:0812.3962.

\providecommand{\CM}{C_{M_{24}}}
\providecommand{\U}{\mathrm{U}(1)}
\providecommand{\Z}{\mathbb{Z}}
\providecommand{\kBKM}{\kappa_{\mathrm{BKM}}}
\providecommand{\kcat}{\kappa_{\mathrm{cat}}}
\providecommand{\kch}{\kappa_{\mathrm{ch}}}
\providecommand{\kfib}{\kappa_{\mathrm{fiber}}}
\providecommand{\aKthree}{[\alpha_{K3}]}
\providecommand{\res}{\operatorname{res}}

\section*{$H^3$ transgression for Mathieu twisted twining}

\paragraph{The Mathieu anomaly class.}
Gaberdiel, Persson, Ronellenfitsch and Volpato construct the
discrete anomaly class of the Mathieu moonshine K3 model as a primitive
class
\[
  \aKthree\in H^3\bigl(M_{24},\U\bigr)\cong\Z/12.
\]
The computation of \(H^3(M_{24},\U)\) is group cohomology in degree
three with \(\U\)-coefficients. The Schur multiplier
\(H^2(M_{24},\U)\) is a separate degree-two object. The exponential
sequence identifies \(H^n(G,\U)\) with \(H^{n+1}(G,\Z)\) for finite
\(G\) and \(n\geq 1\), so a Schur multiplier computation supplies a
different integral degree from the Mathieu anomaly class. The K3
anomaly used by GPV is the generator of the displayed \(\Z/12\),
fixed by the multiplier systems of the twining genera.

\paragraph{Dijkgraaf, Pasquier and Roche transgression.}
Let \(\alpha\) be a normalized representative of a class in
\(H^3(M_{24},\U)\). For \(g\in M_{24}\) and
\(x,y\in \CM(g)\), the DPR slant product is represented by
\[
  c_g(x,y)=
  \frac{\alpha(g,x,y)\alpha(x,y,g)}
       {\alpha(x,g,y)}.
\]
It defines a two-cocycle on \(\CM(g)\), hence a natural map
\[
  \tau_g\colon H^3(M_{24},\U)\longrightarrow H^2(\CM(g),\U),
  \qquad [\alpha]\longmapsto [c_g].
\]
This centralizer \(H^2\)-class governs the projective action of
\(\CM(g)\) on the \(g\)-twisted sector. The cyclic restriction
\[
  \res^{M_{24}}_{\langle g\rangle}\aKthree
  \in H^3(\langle g\rangle,\U)
\]
has a different role: it is the level-matching class of the cyclic
orbifold generated by \(g\).

\paragraph{Centralizer \(H^2\) and cyclic \(H^3\).}
The large-prime Mathieu classes have the following centralizer and
cyclic cohomology data.
\[
\begin{array}{c|c|c|c}
g & \CM(g) & H^2(\CM(g),\U) & H^3(\langle g\rangle,\U)\\
\hline
7A,7B & S_3\times C_7 & 0 & \Z/7\\
11A & C_{11} & 0 & \Z/11\\
23A,23B & C_{23} & 0 & \Z/23
\end{array}
\]
For \(C_N\), the periodic resolution gives
\[
  H^2(C_N,\U)=0,\qquad H^3(C_N,\U)\cong \Z/N.
\]
For \(S_3\times C_7\), the Schur product formula gives
\[
  H^2(S_3\times C_7,\U)\cong
  H^2(S_3,\U)\oplus H^2(C_7,\U)\oplus
  \operatorname{Hom}(S_3^{\mathrm{ab}}\otimes C_7,\U)=0,
\]
because the Schur multipliers of \(S_3\) and \(C_7\) vanish and
\(S_3^{\mathrm{ab}}\cong C_2\) has zero tensor product with \(C_7\).
Thus \(\tau_g(\aKthree)=0\) in centralizer cohomology for the five
displayed classes. A nonzero cyclic value in
\(H^3(\langle g\rangle,\U)\) requires an explicit restriction of the
chosen global three-cocycle. GPV's list of twining genera with
nontrivial multiplier contains the classes
\[
  2B,\ 3B,\ 4A,\ 4C,\ 6B,\ 10A,\ 12A,\ 12B,\ 21A,\ 21B.
\]
The classes \(7A,7B,11A,23A,23B\) are absent from that list, and GPV
records their centralizer representations as linear. The residues
\(3\in\Z/7\), \(2\in\Z/11\), and \(1\in\Z/23\) are therefore
unsupported by GPV and by the first-principles computation above.

\paragraph{Simultaneous \((g,h)\)-twining.}
For a commuting pair with \(h\in\CM(g)\), Persson and Volpato use the
class \([c_g]\in H^2(\CM(g),\U)\) to define the projective
\(\CM(g)\)-representation and the \(c_g\)-regularity condition on
the conjugacy class of \(h\). A statement indexed only by the orders
of \(g\) and \(h\) is underdetermined. The data are the actual
commuting representatives, the subgroup \(\langle g,h\rangle\), the
centralizer class of \(h\), and the restriction of \(\alpha_{K3}\) to
that subgroup. In the cyclic test generated by an element \(x\) of
class \(6A\), with \(g=x^3\) of class \(2A\) and \(h=x^2\) of class
\(3A\), the relevant restriction lies in
\[
  H^3(\langle x\rangle,\U)\cong H^3(C_6,\U)\cong \Z/6.
\]
Geometric Mukai subgroups inside \(M_{23}\) give the trivial
multiplier cases in the GPV reduction, while a general simultaneous
twining statement requires the actual subgroup calculation above.

\paragraph{CHL constants.}
For the CHL levels used in the K3 by elliptic-curve Borcherds
denominator, the constant term and Borcherds weight are
\[
\begin{array}{c|c|c}
N & c_N(0) & \kBKM(\Phi_N)\\
\hline
1 & 10 & 5\\
2 & 8 & 4\\
3 & 6 & 3\\
4 & 4 & 2\\
6 & 2 & 1
\end{array}
\]
with
\[
  \kBKM(\Phi_N)=\frac{c_N(0)}{2}.
\]

\paragraph{The K3 by elliptic-curve invariant split.}
The invariants attached to \(K3\times E\) occupy separate slots.
\[
\begin{array}{c|c}
\text{invariant} & \text{value}\\
\hline
\kcat(K3\times E) & 0=\chi(\mathcal O_{K3})\chi(\mathcal O_E)=2\cdot 0\\
\kch^{\mathrm{compact}}(K3\times E) & 0=1-1+1-1\\
\kch^{\mathrm{Heis}}(K3\times E) & 3\\
\kBKM(\mathfrak g_{\Delta_5}) & 5=c_1(0)/2\\
\kfib(K3) & 24
\end{array}
\]
The total-space categorical invariant, compact Hodge supertrace,
Heisenberg chiral rank, Borcherds weight, and K3 fibre Euler number
are distinct entries.

\paragraph{The Gritsenko and Clery dd catalogue.}
The eight dd-modular Borcherds products form a catalogue indexed by
\((t,N)\) and by weight \(k\).
\[
\begin{array}{c|c|c|c}
\text{form} & (t,N) & k & 2k\\
\hline
M_5(\Gamma_1,\chi_2)=\Delta_5 & (1,1) & 5 & 10\\
M_2(\Gamma_2,\chi_4) & (2,1) & 2 & 4\\
M_3(\Gamma^{(2)}_0(2),\chi_2) & (1,2) & 3 & 6\\
M_1(\Gamma_3,\chi_6) & (3,1) & 1 & 2\\
M_2(\Gamma^{(2)}_0(3),\chi_2) & (1,3) & 2 & 4\\
M_{1/2}(\Gamma_4,\chi_8) & (4,1) & 1/2 & 1\\
M_{3/2}(\Gamma^{(2)}_0(4),\chi_4) & (1,4) & 3/2 & 3\\
M_1(\Gamma_2(2),\chi_4) & (2,2) & 1 & 2
\end{array}
\]
This table is separately indexed from the CHL table. The two tables
meet at the \(\Delta_5\) row, where \(c_1(0)=10\) and
\(\kBKM(\mathfrak g_{\Delta_5})=5\).
